\chapter{Tela de diagnóstico de entrada}
\label{ap:diagnostico}

Este apêndice documenta o instrumento de calibração de controles do projeto,
que corresponde aos requisitos FR-007 (diagnóstico) e FR-008 (conexão e
desconexão a quente). A tela lista os dispositivos conectados, os eixos e os
botões em tempo real, exibe o estado da última conexão ou desconexão e permite
ajustar a zona morta aplicada a todas as ações de \emph{gameplay}. A
Figura~\ref{fig:input_test} registra a tela em execução na máquina de
referência.

\begin{figure}[htbp]
\centering
\includegraphics[width=0.86\textwidth]{input_test.png}
\caption{Tela de diagnóstico de entrada em execução: dispositivos detectados, eixos com valores instantâneos, botões pressionados, zona morta configurável e mensagem de estado da conexão.}
\label{fig:input_test}
\par\vspace{2pt}{\footnotesize Fonte: captura do autor (2026), gerada pela cena de captura do projeto em 1280$\times$720.}
\end{figure}

A tela cumpre dois papéis no projeto. O primeiro é de engenharia: como o
adaptador USB de controle (modelo de PlayStation~1) enumera botões e eixos de
maneira imprevisível, a calibração precisa ser observável --- e é essa tela
que permite descobrir, em segundos, qual índice físico corresponde a cada
ação, sem depender de tentativa e erro no jogo. O segundo papel é de
auditoria: a tela é a única parte do projeto autorizada a ler índices
numéricos de botão e eixo; todo o \emph{gameplay} consome apenas ações
semânticas, e essa separação é verificável no código.
